\chapter{问题、状态与闭环}

应用对象是一台智能蒸烤一体机。操作者放进一块食品，选择蒸、烤或组合，然后关门。机器在门内能改的量只有四个：加热器功率、蒸汽、排汽和风机。表面温度、腔温和质量可以读到。冷点、内部热阻和反应进度读不到。

后面各章都收束到这四个执行量。多尺度模型说明食品内部为什么不能由表面温度代替。估计器说明哪些量能从可安装的传感器里恢复。探针说明一次短时的蒸汽和风是在加工，还是为了看清内部热阻。因果分解说明同一次蒸汽沿着哪几条路径改变这块食品。安全投影说明品质模型出错时，加热器这一路仍由它自己的温度测量接管。

一维十六条控制体是这块食品在腔体里的工作模型。它不是整机的外形，也不代替门、水位和漏电。那些硬测量在安全一章里单独列出，因为它们不读取品质。

本章要判定的量是：一次热加工失败，究竟缺的是传感器清单，还是缺决策泛函的充分统计量。判定结果决定后文哪些观测允许进入向量，哪些边界变化必须离开平衡，以及一次运行缺了哪一步就不得进入比较。

\section{信息赤字}

\paragraph{要判定的量}
在时刻 \(t\)，可安装观测 \(\vect{y}_{0:t}\) 对决策泛函 \(\phi(\vect{x})\) 的剩余不确定度，是否仍然高于允许误判。高于允许误判就称为信息赤字。蒸制内部未到、表面已经过熟、开环时间对不上这一块的厚度，以及干烧、超温、漏电，都要能放进同一个判定，而不是分成“烹饪功能”和“安全功能”两张互不引用的表。

\paragraph{继承的文献层与不继承的层}
不确定度用熵来写，继承的是实验所提供信息的那一层含义：先验熵与后验熵之差 \cite{lindley1956}。不继承把熵本身当作要发明的新物理量。观测方程的病态性继承逆热传导：表面读数的高频误差随厚度被放大 \cite{ozisik2000}。不继承某一热像仪的标称精度。腔体作为边界而食品作为多孔介质，继承体积平均后的热湿传递 \cite{whitaker1977,datta2007}。不继承把“蒸”写成一个更大的导热系数。

\paragraph{方程}
真实状态 \(\vect{x}(t)\in\cX\)，传感器给出 \(\vect{y}(t)=h(\vect{x}(t))+\vect{v}(t)\)。对只含停机、表面品质和安全的泛函 \(\phi\)，
\begin{equation}
\Delta\mathcal{I}(t)=H\big(\phi(\vect{x}(t))\mid \vect{y}_{0:t}\big)-H_{\mathrm{tol}}.
\end{equation}
\(H_{\mathrm{tol}}\) 是该泛函的误判上限，由动作切换的滞回带决定，不由传感器样本方差单独决定。\(\Delta\mathcal{I}>0\) 时若仍按开环时间表改边界，冷点进度和表面干耗会分开漂移：二者对同一边界的符号不必相同。

系统要完成的链是
\begin{equation}
\begin{aligned}
&\text{边界激励}\to\text{多尺度状态}\to\text{后验}\\
&\to\text{信息价值与品质价值}\to\text{安全投影}\to\text{新的边界}.
\end{aligned}
\end{equation}
链上缺一环，后面的环就没有输入。例如没有后验，信息价值和品质价值无法分开；没有安全投影，品质模型的错误会直接变成执行器指令。

\paragraph{本节不能推出的结论}
\(\Delta\mathcal{I}\) 的定义不给出微生物的绝对风险，也不给出某一菜谱的分钟数。\(H_{\mathrm{tol}}\) 尚未由外部温度历程标定。把赤字写成“还要再加一个探针”是产品语言，不是这条定义。

\section{五个必须同时成立的性质}

\paragraph{要判定的量}
下面五条是否同时成立。只成立其中几条时，系统会退回到已经有的装置：只会定时，只会跟腔温，或只会在超温时断电。

\paragraph{继承的文献层与不继承的层}
可重构继承动态系统的可观测性，但工作方程用的是两节点热容，不是一般的线性系统教科书例题。可分离继承对偶控制：输入同时推动状态和辨识 \cite{feldbaum1965}。不继承任意波形的连续最优控制，家电执行器只有有限探针库。可迁移继承“类别只改先验”这一用法，不继承把三种食材写成三套互不相关的方程。可更新继承层次贝叶斯里精度加权，不继承用一条轨迹覆盖整个类别。可拦截继承控制屏障：安全集由测量定义，参考动作只进入目标函数 \cite{ames2017cbf}。不继承把屏障函数的定义本身算作新对象。

\paragraph{五条性质}
\begin{enumerate}
\item \textbf{可重构。} 冷点温度、表面含水、反应进度和内部热阻，能由有限观测与短时边界激励更新。稳态观测不含内部热阻，所以“可重构”依赖离开平衡的激励，不是依赖更密的腔温采样。
\item \textbf{可分离。} 同一次阀门动作，要么为了降低后验熵，要么为了推动品质泛函。目标函数不同，日志里要能查出这一时刻是哪一个函数发出的指令。
\item \textbf{可迁移。} 宏观传热、孔隙水分和反应进度耦合成低维决策状态。食材类别只改变先验和终端目标，不改变状态的分量定义。
\item \textbf{可更新。} 单次过程只更新事先声明的传递参数。结构方程的父节点不随这次数据改写。
\item \textbf{可拦截。} 品质模型的输出在进入执行器之前，被只依赖硬测量的不等式投影。硬保护不经过该模型，也不经过产生参考动作的处理器。
\end{enumerate}

\paragraph{与已运行核的距离}
已运行的一维核实现了可分离的一部分：对偶策略在固定时间窗里改蒸汽和风机，其余时间用冷点估计决定加热，停机后再降功率。它没有实现一般的可重构：热阻更新使用了真值温差做成的参照，见附录 \ref{app:discrete}。可拦截只落实到加热器温度这一路，而且满功率平衡在 \(196.5^\circ\mathrm{C}\)，\(230^\circ\mathrm{C}\) 削减不会动作。因此五条是设计判据，不是这次日志已经逐条定量成立的清单。

\paragraph{本节不能推出的结论}
五条同时成立，不能推出某一策略的中位误差更小。误差比较放在第 \ref{ch:sim} 章，而且只引用已经落盘的列。

\section{决策量}

\paragraph{要判定的量}
操作者需要的判断只有四个：能否停止，表面含水是否偏离目标，温度是否还在散开，继续施加边界是否安全。每个判断对应哪些隐藏量，决定观测向量的最小组成。

\paragraph{继承的文献层与不继承的层}
“熟了不是到达某一温度”继承食品品质动力学：速率是温度的指数，进度是速率的积分 \cite{vanboekel2008}。热力当量继承等效时间积分 \cite{holdsworth1997}。不继承把 \(0.85\) 写成认证限值。安全判断继承家用电器对漏电、温升和异常工作的通用要求 \cite{iec60335,gb4706}，只作为测量点必须存在的入口，不作为本文已经满足条款的声明。

\begin{table}[ht]
\centering
\caption{四个判断及所需隐藏量}
\begin{tabularx}{\textwidth}{lXX}
\toprule
判断 & 隐藏量 & 信息来源 \\
\midrule
能否停止 & 冷点进度的下分位 & 逆估计与探针确认 \\
表面是否偏离 & 表面含水、质量损失 & 质量、湿度、冷凝 \\
是否仍不均匀 & 温度离散度 & 热像与降阶场 \\
继续是否安全 & 干烧、漏电、腔压、门 & 不进入品质模型的硬测量 \\
\bottomrule
\end{tabularx}
\end{table}

腔温只描述空气。新增一路观测，必须降低上表中至少一个判断的 \(\Delta\mathcal{I}\)，否则不进入观测向量。穿刺温度能降低“能否停止”的不确定度，但它改变局部场，而且不刺中冷点时读数不代表 \(\operatorname{ess\,inf} T\)。因此穿刺只作为对照通道，不进入主观测向量，系统的可重构性不依赖它。

\paragraph{本节不能推出的结论}
四个判断没有覆盖风味分子和微生物初始数量。那些量不在状态里。表中“信息来源”是规格上的通道；本次核只有表面温度、平均含水和腔温的带噪读数，热像、质量和漏电没有采样。

\section{闭环的完整性条件}

\paragraph{要判定的量}
一次运行能否进入结果比较。判据是六步是否都留下记录，而不是最终误差是否好看。

\paragraph{继承的文献层与不继承的层}
继承的是估计与控制分离：真值只用于打分 \cite{rawlings2017}。不继承把同一套方程自洽积分的残差叫做感知误差。

\begin{enumerate}
\item 真值模型生成食品与腔体边界，参数对在线模块不可见。
\item 在线模块只接收可安装通道的带噪采样。
\item 后验给出决策状态、信息赤字和建议执行量。
\item 执行量经安全投影写回真值边界。
\item 评分使用真值场上的进度、均匀性、能耗和越界次数。
\item 探针前后的协方差若未下降，该次动作记为无效信息，不得改变停止时刻。
\end{enumerate}

真值写入在线模块会使第 2 步失效，该次数据作废。已运行核的热阻更新在“是否更接近参照”这一步使用了由真值温差构成的 \(r_{\mathrm{proxy}}\)。这一步不满足第 2 步的严格读法。正文报告方差下降时，必须同时报告这个参照，不能把下降写成纯观测更新。

第 6 步在核里的实现是：对偶策略在厚件上，若方差仍大于 \(0.45\) 且厚度大于 \(40\,\mathrm{mm}\)，否决停止。方差乘 \(0.35\) 只在新值更接近参照时发生。无效探针不能移动停止时刻，这条在规则上存在；它不是 Fisher 矩阵的特征值检验。

\paragraph{本节不能推出的结论}
六步写全，不能推出主结论已经由异构真值模型支持。本次积分器是十六条控制体，不是带压力方程的多孔介质真值。同构自检不能放进主表。

\section{文档结构}

后文按这条链展开。链的两端都是这台蒸烤一体机：入口是四个边界，出口仍是四个边界。中间各章不另设一套执行器。

\begin{center}
\begin{tabularx}{\textwidth}{llX}
\toprule
章 & 对象 & 在闭环中的位置 \\
\midrule
2 & 已有模型与装置能力 & 划定新增环节 \\
3 & 层、时间尺度、模块合同 & 规定信息不能走捷径 \\
4 & 多尺度场与决策状态 & 生成被估计的对象 \\
5 & 观测、病态性、滤波 & 被动信息的上限 \\
6 & 对偶控制与探针 & 主动信息 \\
7 & 结构因果与路径 & 解释与干预 \\
8 & 算子与参数更新 & 跨时间的模型 \\
9 & 品质优化与安全投影 & 执行 \\
10 & 真值隔离的试验 & 闭环是否成立 \\
11 & 可实施步骤与实施例 & 把公式收成腔体上的步骤 \\
12 & 命题、证明与图件 & 可否定的陈述 \\
13 & 交互演示与数据发布 & 数值核、界面与下载字段 \\
14 & 文献分层 & 每条方程继承与不继承的结果 \\
15 & 图件总表 & 88 个图位的坐标和归属 \\
16 & 成稿 & 同一正文的 PDF 与 Word \\
17 & 外部数据 & 哪些参数被公开数据约束 \\
18 & 构建顺序 & 依赖与尚未执行的规格 \\
\bottomrule
\end{tabularx}
\end{center}

每一章内部固定四段：要判定的量，继承与不继承的层，方程或协议，不能推出的结论。摘要中的定量句必须指向已经画出的数据图。没有数据图的陈述只用条件句，并放在定义或命题里。
